\section{A common collar and weighted monotone convergence}
\label{sec:monotone}

The results of this section use the internal inputs listed in
Section~\ref{sec:inputs}. In particular, they do not assert an independent
proof of the regularization theorems in \textup{R} and \textup{B}.
All currents and test forms belong to the original CLD calculus.
Positivity means positivity in the cone used in \textup{R}, not strong
positivity. Set $D=d'd''$.

We first record the exact local consequence of the internal approximation
results that is needed below. A smooth psh function in this section has a
positive original $(1,1)$ curvature \emph{form}.

\begin{lemma}[A common domain for smooth majorants]
\label{lem:common-majorants}
Let $x\in U$, and let $u_{i,j}\downarrow u_i$ be finitely many decreasing
families in $\Cb(U)$ with locally bounded limits. After shrinking around
$x$, there is one ordinary neighborhood $W$, independent of $i,j$, with
the following property. Each of the functions $a=u_{i,j}$ and $a=u_i$
has smooth psh majorants $a_\nu$ on $W$ such that
\[
 a_\nu\geq a,\qquad a_\nu(x')\longrightarrow a(x')\quad(x'\in W).
\]
The majorants are bounded uniformly in $\nu$ on each fixed compact subset
after a finite prefix is discarded. If $a\leq M$ on a compact neighborhood
$C\Subset W$, its majorants may be chosen with $a_\nu\leq M+1$ on $C$.
The index discarded may depend on the function $a$.
\end{lemma}

\begin{proof}
Use \textup{B}, Proposition 7.1, to work on one ordinary algebraic
neighborhood, with the original field and calibration. Apply \textup{B},
Lemma D.2, to the entire finite family of sequences, before choosing
approximants. It gives one projective variety, one ample line, one frame
on $W$, and bounded classical metrics whose weights on $W$ are the given
functions plus one fixed constant. By \textup{B}, Theorem 4.4, every such
metric is a decreasing limit of pure Fubini--Study metrics on this same
line. Their weights $p_\nu$ are continuous, satisfy $p_\nu\geq a$, and
decrease pointwise to $a$ after the fixed constant is removed.

Smooth each pure Fubini--Study metric uniformly, using \textup{R},
Lemma 3.2 and Proposition A.1. If the uniform error is at most
$\epsilon_\nu\downarrow0$, add $\epsilon_\nu$ to the smoothed weight.
The resulting smooth psh function $a_\nu$ satisfies
\[
 p_\nu\leq a_\nu\leq p_\nu+2\epsilon_\nu.
\]
This proves pointwise convergence and the majorant property.
On $C$, the increasing open sets $\{p_\nu<M+1/2\}$ cover $C$ whenever
$a\leq M$. Compactness gives a common index on $C$. Discard this prefix
and require $2\epsilon_\nu<1/2$. The same argument with any bound for
$a$ gives local uniform boundedness of the row. No monotonicity of the
\emph{smoothed} row is asserted or needed.
\end{proof}

\begin{lemma}[A smooth psh collar]
\label{lem:smooth-collar}
Under the assumptions of Lemma~\ref{lem:common-majorants}, there are
compact polynomial norm boxes
\[
 x\in\operatorname{Int}C_1,\qquad
 C_1\subset\operatorname{Int}C_{\mathrm{mid}},\qquad
 C_{\mathrm{mid}}\subset\operatorname{Int}C_2\Subset W,
\]
a smooth psh function $b$ on a neighborhood of $C_2$, and a nonnegative
original smooth cutoff $\chi$ compactly supported in
$\Omega=\operatorname{Int}C_2$, with the following properties.
After adding one constant $c$ to all the original functions, put
\[
 v_{i,j}=\max\{u_{i,j}+c,b\},\qquad
 v_i=\max\{u_i+c,b\}.
\]
Then:
\begin{enumerate}
\item $v_{i,j},v_i\in\Cb(\Omega)$ and $v_{i,j}\downarrow v_i$;
\item $v_{i,j}=u_{i,j}+c$ and $v_i=u_i+c$ on a neighborhood of $C_1$;
\item $\chi=1$ on a neighborhood of $C_{\mathrm{mid}}$;
\item all $v_{i,j},v_i$ equal $b$ on one fixed neighborhood of
      $\operatorname{supp}d'\chi\cup\operatorname{supp}d''\chi$;
\item each patched function $v$ has smooth psh majorants $v_\nu\geq v$
      converging pointwise to $v$, uniformly bounded on $\operatorname{supp}\chi$,
      and equal to $b$ on that same collar.
\end{enumerate}
The majorant rows in (5) need not be decreasing.
\end{lemma}

\begin{proof}
Choose the boxes by the construction in \textup{B}, Lemma D.2. In
particular, lower-bound polynomials have no zero near $C_2$. Choose
thresholds strictly between the defining thresholds for $C_1$ and
$C_{\mathrm{mid}}$. The finite maximum
\[
 b_0=\max\{0,\log|P_a^{N_a}/c_a|,
                     \log|d_l/Q_l^{M_l}|\}
\]
is zero on a neighborhood of $C_1$ and at least some $\delta>0$ on a
neighborhood in $C_2$ of $C_2\setminus\operatorname{Int}C_{\mathrm{mid}}$.
All its branches other than the zero branch are logarithms of analytic
units where they are active. A branch with a zero is strictly inactive
near that zero.

Apply the regularized maximum of \textup{R}, Proposition A.1, with
smoothing radius smaller than all the fixed gaps. The result $\widetilde b$
is an original smooth psh function: zero branches at polynomial zeros
are deleted by the exact inactivity property. On a neighborhood of
$C_1$, only the constant zero branch is active, so $\widetilde b=0$.
On the outer collar, $\widetilde b\geq\delta$.

The endpoint bounds give a constant $c$ and a finite $M$ such that
\[
 1\leq u_i+c\leq u_{i,j}+c\leq M
 \quad\hbox{on a neighborhood of }C_2.
\]
Multiply $\widetilde b$ by a positive constant so that the resulting $b$
is greater than $M+3$ on the outer collar. Choose $\chi$, with
$0\leq\chi\leq1$, equal to one near $C_{\mathrm{mid}}$ and compactly
supported in $\Omega$. Smooth cutoff existence and a finite localization
are part of the original CLD test calculus. Its derivative support lies
where the fixed outer-collar inequality holds.

Classical psh functions are stable under maxima and addition of constants.
This proves (1)--(4). To prove (5), apply Lemma~\ref{lem:common-majorants}
to the \emph{unpatched} function $a=u_{i,j}+c$ or $u_i+c$. Discard a
prefix so that its smooth psh majorants $a_\nu$ are at most $M+1$ on
$C_2$. Define
\[
 v_\nu=\operatorname{rmax}_{\epsilon_\nu}(a_\nu,b),
 \qquad \epsilon_\nu\downarrow0,
\]
with all radii less than the fixed collar gap. The regularized maximum
majorizes the ordinary maximum and has error at most $\epsilon_\nu$.
It is exactly $b$ on the collar because the other branch is inactive
there. These facts give all the assertions in (5).
\end{proof}

\begin{remark}
Neither lemma supplies a strictly psh exhaustion of an arbitrary analytic
domain. The collar is constructed on a chosen small norm box. The closed
positive form used to detect a test below is chosen separately and may
be degenerate in directions not occurring in that test.
\end{remark}

\begin{lemma}[Green exchange on a common collar]
\label{lem:collar-green}
Let $Q$ be a symmetric positive $d'$- and $d''$-closed current of
bidegree $(n-1,n-1)$ on $\Omega$. Let $\chi\in A_c^{0,0}(\Omega)$.
Suppose that $a$ is smooth, $g$ is locally bounded Borel, and both equal
one smooth function $b$ on a neighborhood of
$\operatorname{supp}d'\chi\cup\operatorname{supp}d''\chi$.
Whenever $D(gQ)$ has coefficient Radon type, one has
\[
 \langle aD(gQ),\chi\rangle
 =\langle gD(aQ),\chi\rangle.
\]
In this statement positivity of $a$ and $g$ is not required.
\end{lemma}

\begin{proof}
Since $Q$ is closed and $a$ is smooth, $D(aQ)=Da\wedge Q$. The
definition of the distributional derivative gives
\begin{align*}
 \langle aD(gQ)-gD(aQ),\chi\rangle
  &=\langle gQ,D(\chi a)-\chi Da\rangle\\
  &=\langle Q,g\{aD\chi+d'\chi\wedge d''a
                         +d'a\wedge d''\chi\}\rangle.
\end{align*}
Every term in braces is supported in the collar. Substituting $a=g=b$
there makes the last expression
\[
 \langle bQ,D(\chi b)-\chi Db\rangle=0.
\]
The final equality follows by applying the same distributional identity
to the smooth pair $(b,b)$. The substitutions are valid identities of
coefficient measures. No derivative of the Borel function $g$ is used.
\end{proof}

\begin{lemma}[Monotonicity of the ordered trace integral]
\label{lem:trace-monotonicity}
Let $S$ have bidegree $(p,p)$, let $0\leq k\leq n-p$, and let
$\alpha$ be a smooth, positive, $d'$- and $d''$-closed form of bidegree
$(n-p-k,n-p-k)$ such that wedging it with a positive current preserves
positivity. Let $\chi\geq0$ be compactly supported in $\Omega$.

For $0\leq i\leq k$, suppose $f_i,g_i\in\Cb(\Omega)$,
$f_i\leq g_i$, and all these functions equal the same smooth psh
function $b$ near the derivative support of $\chi$. Assume that each
$g_i$ has the smooth majorants with this exact collar supplied by
Lemma~\ref{lem:smooth-collar}. Define, using the ordered products,
\[
 J(f_0;f_1,\ldots,f_k)
   =\langle f_0T_{S,k}(f_1,\ldots,f_k),\chi\alpha\rangle.
\]
Then
\[
 J(f_0;f_1,\ldots,f_k)
 \leq J(g_0;g_1,\ldots,g_k).
\]
Only existence and positivity of the ordered products are used here.
\end{lemma}

\begin{proof}
The case $k=0$ is measure monotonicity. For $k>0$, fix smooth psh
majorants $G_i\geq g_i$ with the common collar. First replace the weight
$f_0$ by $G_0$, using positivity of
$T_{S,k}(f_1,\ldots,f_k)\wedge\alpha$. Apply
Lemma~\ref{lem:collar-green} with
\[
 Q=T_{S,k-1}(f_2,\ldots,f_k)\wedge\alpha
\]
to move $G_0$ into a differentiated position and $f_1$ into the weight.
The resulting current $DG_0\wedge Q$ is positive and closed. Replace
the weight $f_1$ by $G_1$, and repeat. At the end one obtains
\[
 J(f_0;f_1,\ldots,f_k)
 \leq\langle G_k\,DG_0\wedge\cdots\wedge DG_{k-1}\wedge S,
                     \chi\alpha\rangle.
\]
A final application of the same Green identity, together with
commutation of smooth even-degree factors, makes the right-hand side
$J(G_0;G_1,\ldots,G_k)$.

Now let the index of $G_k$ tend to infinity, keeping every other index
fixed. Bounded convergence against the fixed coefficient measures of
$S$ implies $D(G_kS)\to D(g_kS)$. The outer factors and the weight are
smooth and fixed, so this convergence can be tested directly. Next let
the index of $G_{k-1}$ tend to infinity. The background
$D(g_kS)$ is now fixed, so the same argument applies. Continue in the
order $k,k-1,\ldots,1$. Finally, with the ordered product fixed, let the
weight $G_0$ tend to $g_0$ by bounded convergence. The result is exactly
$J(g_0;g_1,\ldots,g_k)$ in its original order. At no step are two
discontinuous limits interchanged.
\end{proof}

\begin{lemma}[An upper bound for decreasing weights]
\label{lem:upper-weight}
Let $P_j,P$ be positive coefficient Radon currents of one bidegree,
with $P_j\to P$ on all original compact smooth tests. Let $a_j\downarrow a$
be locally uniformly bounded functions. Suppose that, near the support
of a fixed positive compact test $\Theta$, every fixed $a_l$ has bounded
continuous majorants converging pointwise to $a_l$.
Then
\[
 \limsup_{j\to\infty}\langle a_jP_j,\Theta\rangle
       \leq\langle aP,\Theta\rangle.
\]
\end{lemma}

\begin{proof}
For $j\geq l$ and a continuous majorant $a_{l,\nu}\geq a_l$,
positivity gives
\[
 \langle a_jP_j,\Theta\rangle
       \leq\langle a_{l,\nu}P_j,\Theta\rangle.
\]
Keep $l,\nu$ fixed and pass to the limit in $j$, using \textup{R},
Lemma 3.1. Bounded convergence against the fixed positive Radon measure
$\mu_{P,\Theta}$ first allows $\nu\to\infty$ and then $l\to\infty$.
\end{proof}

\begin{theorem}[Weighted decreasing continuity for a fixed arbitrary background]
\label{thm:relative-monotone}
Assume the internal inputs \textup{R} and \textup{B}. Let $S$ be any
symmetric positive original CLD current of bidegree $(p,p)$ on $U$,
separately $d'$- and $d''$-closed. Let $0\leq k\leq n-p$.
If $u_{i,j}\downarrow u_i$ in $\Cb(U)$ for $0\leq i\leq k$, with
locally bounded limits, then
\[
 u_{0,j}T_{S,k}(u_{1,j},\ldots,u_{k,j})
 \longrightarrow
 u_0T_{S,k}(u_1,\ldots,u_k)
\]
on every original compact smooth complementary test form. The ordered
unweighted products have the corresponding joint decreasing continuity.
The background $S$ is fixed throughout.
\end{theorem}

\begin{proof}
Write $W_k$ for the weighted assertion with $k$ differentiated slots.
Bounded convergence for the fixed coefficient measures of $S$ proves
$W_0$. Distributional differentiation shows that $W_{k-1}$ implies
joint unweighted convergence in degree $k$. We prove $W_k$ assuming
this unweighted convergence.

The statement is local on compact test supports. Use a finite original
smooth partition to reduce to one test $\eta$ supported in a small
norm box $C_1$ as in Lemma~\ref{lem:smooth-collar}. Include in the chosen
ordinary neighborhood all the unit coordinates needed to present
$\eta$. Let these coordinates be $t_1,\ldots,t_N$, and set
\[
 \beta=\sum_{a=1}^N d't_a\wedge d''t_a,
 \qquad r=n-p-k,\qquad \alpha=\frac{\beta^r}{r!}.
\]
For $r=0$ set $\alpha=1$. The form $\alpha$ is closed and strongly
positive. In the fixed finite coefficient space it is a positive
definite trace form. After replacing $\eta$ by its symmetric part,
there is $C>0$ such that
\[
 \Theta_\pm=C\chi\alpha\pm\eta
\]
are positive original compact tests. This is precisely the matrix
domination used in \textup{R}, Lemma 3.1; additional scalar coordinates
are handled by pullback to a common refinement. No definiteness in
all possible moment directions is required.

Patch all the potentials to $v_{i,j},v_i$ by
Lemma~\ref{lem:smooth-collar}, and put
\[
 P_j=T_{S,k}(v_{1,j},\ldots,v_{k,j}),\qquad
 P=T_{S,k}(v_1,\ldots,v_k).
\]
The induction hypothesis gives $P_j\to P$. Applying
Lemma~\ref{lem:trace-monotonicity} to the ordered tuples
$(v_0,\ldots,v_k)\leq(v_{0,j},\ldots,v_{k,j})$ yields
\[
 \langle v_0P,\chi\alpha\rangle
       \leq\langle v_{0,j}P_j,\chi\alpha\rangle.
\]
Lemma~\ref{lem:upper-weight} gives the reverse limsup inequality.
Consequently
\begin{equation}
 \label{eq:trace-limit}
 \langle v_{0,j}P_j,\chi\alpha\rangle
       \longrightarrow\langle v_0P,\chi\alpha\rangle.
\end{equation}
Apply Lemma~\ref{lem:upper-weight} to $\Theta_+$ and $\Theta_-$.
Together with \eqref{eq:trace-limit}, the two resulting inequalities
give respectively the upper and lower bounds for
$\langle v_{0,j}P_j,\eta\rangle$. Thus this pairing converges to
$\langle v_0P,\eta\rangle$.

On a neighborhood of $\operatorname{supp}\eta$, the patched potentials
are $u_{i,j}+c$ and $u_i+c$. Ordinary locality of the recursion identifies
the differentiated products there with the original products, since
constants in differentiated slots vanish. Remove the constant $c$ from
the weight by the unweighted degree-$k$ convergence already established.
Finite summation over the test partition proves $W_k$.

This induction proves all $W_k$ and all allowed unweighted assertions.
It uses neither capacity nor energy convergence, nor any interchange
of two discontinuous approximation parameters.
\end{proof}

\begin{corollary}[Symmetry and the first-stage calculus]
\label{cor:full-calculus}
Under the hypotheses of Theorem~\ref{thm:relative-monotone}, the currents
$T_{S,k}$ are symmetric in the differentiated arguments,
positive, symmetric as currents, and separately $d'$- and $d''$-closed.
They are additive and positively homogeneous in each potential, commute
with ordinary restriction and gluing, and are independent of the
regularization used to compute them. In particular,
\[
 D\bigl(uD(vS)\bigr)=D\bigl(vD(uS)\bigr)
 \qquad(u,v\in\Cb(U),\ p\leq n-2).
\]
\end{corollary}

\begin{proof}
All assertions except permutation symmetry were proved for the ordered
recursion. Choose, locally on one fixed domain, continuous semipositive
PL approximants decreasing to each potential, as in \textup{B}.
Continuous arbitrary-background products are symmetric, by uniform
smooth approximation and fixed-continuous-multiplier convergence.
Theorem~\ref{thm:relative-monotone}, with a constant weight, passes this
identity to the given potentials. The recursion identifies the result
independently of the approximants and on overlaps.
\end{proof}

\begin{corollary}[Top-degree convergence as measures]
\label{cor:top-measures}
If $p+k=n$, the weighted convergence in
Theorem~\ref{thm:relative-monotone} holds against every compactly
supported continuous scalar function, locally as signed Radon measures.
\end{corollary}

\begin{proof}
Unweighted convergence against one nonnegative smooth cutoff equal to
one near a prescribed compact set bounds the masses of the positive
measures there. The weights have a common actual pointwise bound on a
slightly larger compact neighborhood. Thus the weighted variations are
uniformly bounded. Uniform density of original smooth scalar functions,
with a fixed cutoff, extends convergence to continuous scalar tests.
\end{proof}

\begin{remark}[What independence means here]
The construction is intrinsic to the \emph{given potentials} and to $S$.
Independence of regularization does not assert that $D(uS)$ depends only
on the ambient current $D[u]$ for every possible pair of potentials.
That stronger quotient-by-potentials statement is distinct from the
claims proved here and from the special Bedford--Taylor-current
condition in CLD 2025, Definition 8.6.2. Equality on all original tests
also does not assert equality on all $\delta$-tests.
\end{remark}
